\scanpage{11}

\uline{Proposition \unsure{1.11|1.10}.} Let $\mathcal E,\mathcal E'$ be two multi-Galois topoi, $\Pi_1(\mathcal E),\Pi_1(\mathcal E')$ their fundamental groupoids. The obvious functor
\[
\underline{\mathrm{Hom}}_{\mathrm{top}}(\mathcal E,\mathcal E')
\longrightarrow
\underline{\mathrm{Hom}}\bigl(
\underset{C}{\underline{\Pi_1(\mathcal E)}},
\underset{C'}{\underline{\Pi_1(\mathcal E')}}\bigr)
\]
is an equivalence of categories, one finds a quasi-inverse equivalence by composing
\[
\underline{\mathrm{Hom}}(\deleted{\illegible{1}}C,C')
\longrightarrow
\begin{array}[t]{c}
\underline{\mathrm{Hom}}_{\mathrm{top}}(\widehat C,\widehat {C'})\\
\cap\\
\underline{\mathrm{Hom}}(\widehat {C'},\widehat C)
\end{array}
\simeq
\underline{\mathrm{Hom}}_{\mathrm{top}}(\mathcal E,\mathcal E').
\]
\inserted{(\unsure{or|for} $\widehat C\simeq\mathcal E$, $\widehat {C'}\simeq\mathcal E'$)}

\deleted{\illegible{1}}
\uline{1.12} Making explicit the case where $\mathcal E\ \mathcal E'$ are $0$-connected and pointed, hence given as $\mathcal E\simeq\mathrm{Ens}(G)$, $\mathcal E'\simeq\mathrm{Ens}(G')$

\dots

\margin{10}

\scanpage{12}
\mspage{4}

2. \unsure{Application|Applications} to coverings of topoi.

\uline{Theorem 2.1.} Let $\mathcal E$ be a locally connected topos (i.e. every object of which is a sum of $0$-connected objects) and locally simply connected (i.e. admitting a system of generators \unsure{which} are $1$-connected).
\margin{perhaps one must \unsure{suppose|supposed} that $\mathcal E$ has ``enough points'' i.e. enough fiber functors\dots}
\deleted{\illegible{6}}\inserted{Then the category $\mathcal E_{lc}$} of local\textsuperscript{ly} constant objects of $\mathcal E$ is \inserted{a topos} multi-Galois\deleted{\unsure{ne}}, and the inclusion
\[
(2.1.1.)\qquad \mathcal E_{lc}\hookrightarrow\mathcal E
\]
commutes with finite $\varprojlim$ (N.B. in fact, without hypotheses on the topos $\mathcal E$, $\mathcal E_{lc}$ is stable under finite $\varprojlim$) \uuline{and} with arbitrary $\varinjlim$.

\uline{Definition 2.2} We denote this topos by $\mathcal E_{[1]}$, call it the multi-Galois envelope of $\mathcal E$, and the morphism of topoi transposed from the inclusion (2.1.1):

\margin{11}

\scanpage{13}

\[
\mathcal E\longrightarrow\mathcal E_{[1]}
\]
is given the name \uline{canonical morphism}.

\uline{N.B.} This is the equivalent in topos theory of the operation of ``killing the $\Pi_i$ for $i\geq 2$''. One can also define \inserted{an} \deleted{\unsure{$\mathcal E_\infty$}} $\mathcal E_{[0]}$ and a sequence
\[
\mathcal E\unsure{\longrightarrow|-}\mathcal E_{[1]}\unsure{\longrightarrow|-}\mathcal E_{[0]}
\]
[$\mathcal E_0$ is the \uline{discrete} topos defined by $\Pi_0(\mathcal E)$, which has a satisfactory meaning as soon as $\mathcal E$ locally connected\dots]. With suitable hypotheses on $\mathcal E$ (of the type ``local contractibility''), one should be able to define the $\mathcal E_{[i]}$ for $\forall\ i\in\mathbb N$, and canonical morphisms
\[
\mathcal E\longrightarrow\cdots\quad
\mathcal E_{[i]}\unsure{\longrightarrow|-}\mathcal E_{[i-1]}\longrightarrow\cdots\quad
\mathcal E_{[1]}\longrightarrow\mathcal E_0.
\]

\margin{12}

\scanpage{14}
\mspage{5}

2.3. The \uline{fundamental groupoid} of $\mathcal E$ is defined as having for objects the points of $\mathcal E$ (which induce points of $\mathcal E_{[1]}$ thanks to $\mathcal E_{[1]}\longrightarrow\mathcal E$), and for morphisms the \deleted{\unsure{points}} morphisms \uline{of points of} $\mathcal E_{[1]}$. We thus have canonical functors
\[
\underline{\mathrm{Pt}}(\mathcal E)
\xrightarrow{\alpha}\Pi_1(\mathcal E)
\overset{\beta}{\overset{\approx}{\longrightarrow}}
\mathrm{Pt}(\mathcal E_{[1]})
\overset{\text{definition}}{=}\Pi_1(\mathcal E_{[1]})
\]
where $\beta$ is an equivalence (but not surjective on objects), and where of course $\alpha$ is not necessarily an equivalence, nor even fully faithful, or merely faithful. For example if $\mathcal E$ is $1$-connected i.e. $\Pi_1(\mathcal E)$ equivalent to the one-point category, it does not necessarily follow that the Hom \unsure{in|of} $\mathrm{Pt}(\mathcal E)$ are all of cardinality $\leq 1$!

\margin{13}

\scanpage{15}

Since a poin\textsuperscript{t} $P$ of $\mathcal E$ defines a poin\textsuperscript{t} (still denoted $P$ by abuse) of $\mathcal E_{[1]}$, one can therefore define the \uline{universal covering of $\mathcal E$} based at this point, as an object \inserted{$S$} $1$-connected in $\mathcal E_{[1]}$ -- it is characterized \inserted{in $\mathcal E$} by being locally constant, $1$-connected, and endowed with a lifting
\[
P\longrightarrow\mathcal E_{/S}.
\]

But since $\alpha$ is not an equivalence of categories (although it is \deleted{\illegible{1}} \unsure{essentially} surjective if \unsure{one} assumes that $\mathcal E$ has enough points\dots) one can \uline{not} say that $\forall$ universal covering of $\mathcal E$ is defined by a poin\textsuperscript{t} of $\mathcal E$, defined up to unique isomorphism\dots

\margin{14}

\scanpage{16}
\mspage{6}

3. ``pro-multi-Galois'' variants, resp. profinite (restricting ourselves, to simplify, to the case of locally connected topoi\dots).

4. Supplement-remorse on multi-Galois categories, which makes precise the intention that for a topos $\mathcal E$, the datum of an object $S\in\mathcal E$ defines an induced topos $\mathcal E_{/S}\longrightarrow\mathcal E$, and that $S$ is reconstructed up to unique isomorphism from knowledge of the induced topos as a topos \uline{over} $\mathcal E$. Here, \deleted{\illegible{1} $\mathcal E_{/S}$ \unsure{is}} \inserted{$\mathcal E$ being} multi-Galois, $\mathcal E_{/S}$ too -- and the question arises when a morphism of \inserted{\unsure{multi-Galois|multi-Galois}} topoi $\mathcal E'\longrightarrow\mathcal E$ can be regarded as

\margin{15}

\scanpage{17}

an induction morphism. Now if $C=\Pi_1(\mathcal E)$, $C'=\Pi_1(\mathcal E')$, the datum of $\mathcal E'\unsure{\longrightarrow|-}\mathcal E\deleted{\unsure{'}}$ is equivalent to the datum of a functor $C'\longrightarrow C\deleted{'}$. One finds that $\mathcal E'\unsure{\longrightarrow|-}\mathcal E\deleted{\unsure{'}}$ is an induction morphism \unsure{iff} $C'\longrightarrow C\deleted{'}$ is \uline{faithful}. Thus, one finds an equivalence between the category \inserted{$\mathcal E$} (of objects $S$ of the multi-Galois category $\mathcal E\simeq\widehat C$, \deleted{\illegible{1}} where $C$ is a groupoid, \uline{arbitrary if one insists})
\margin{a bit hasty!}
and the category whose objects are the ``groupoids $C'$ over $C$'', with a \uline{faithful} structural functor $C'\longrightarrow C$, the morphisms from $C'_1$ to $C'_2$ being the \inserted{\uuline{isomorphism classes \unsure{of a|of a}}} \deleted{\illegible{1}} pairs $(f,\alpha)$ of a functor $f:C'_1\unsure{\longrightarrow|-}C'_2$ and an

\margin{16}

\scanpage{18}
\mspage{7}

isomorphism of functors $\alpha:p_2f\simeq p_1$
\diagram{%
\begin{tikzcd}
C'_1 \arrow[rr,no head,"f"] \arrow[dr,"p_1"',""{name=pone}]
  & & C'_2 \arrow[dl,"p_2",""{name=ptwo}] \\
  & C &
\arrow[from=pone,to=ptwo,bend left=25,"\overset{\alpha}{\sim}"]
\end{tikzcd}}
\margin{Specify the isomorphisms between pairs $(f,\alpha)$ and \deleted{\unsure{pairs}} $(g,\beta)$\dots}

In the case where for example $C$ is the category reduced to a single object, an automorphism group $G$, this description of the category $\mathcal E=\mathrm{Ens}(G)$ is obviously a bit cumbersome, but it fits well into certain contexts below.

Thus, if $k$ is a base field, the \deleted{\unsure{$\mathcal E$}} category $\mathcal E$ of étale schemes over $k$ is described, in terms of a separable closure \unsure{$k$|the} $k_s$ of $k$ and of the profinite group $\Gamma=\mathrm{Gal}(k_s/k)$, as the \inserted{profinite} groupoids over the profinite groupoid $(\mathrm{pt},\Gamma)$\dots We want to insert this description into a ``Galois'' ``description'' \unsure{of the|of} \inserted{[certain]} schemes [smooth \inserted{quasi-}projective of dimension $\leq 1$] over $k$, at least if $k$ a field of finite type over $\mathbb Q$.

\margin{17}

\scanpage{19}

5. Let $K$ be an extension, of finite type over $\mathbb Q$, and let us choose an algebraic closure $\overline K$ of $K$.
\margin{We set $\Gamma=\mathrm{Gal}(\overline K/K)$}
\inserted{5.1} We consider pairs $(X,S)$, where: \deleted{\unsure{$X$}}

a) $X$ is a projective and smooth scheme over $K$, of dimension $\leq 1$;

b) $S$ is a finite \inserted{reduced} subscheme of $X$, \inserted{(then \illegible{1} \unsure{étale} over $K$)}, contained \deleted{\illegible{2}} \inserted{(in)} \inserted{the union} \unsure{of the} irreducible components of dimension $1$ of $X$.
\margin{(thus finite étale over $\unsure{k|K}$)}

The morphisms $(X',S')\unsure{\longrightarrow|-}(X,S)$ will by definition be the morphisms of schemes
\[
f:X'\longrightarrow X
\]
such that
\[
S'=f^{-1}(S)_{\mathrm{red}}
\]
i.e. such that $\mathrm{supp}\,S'=f^{-1}(\mathrm{supp}\,S)$.

We are looking for a ``Galois description'' of this category, or at least of a full \inserted{sub-}category \inserted{$\mathcal V_K$} which we are going to

\margin{18}

\scanpage{20}
\mspage{8}

define now.

\uline{Lemma} 5.2. Let $\Omega$ be an algebraically closed field, $X$ a \deleted{scheme} \inserted{smooth} projective algebraic curve, connected over $\Omega$, $S$ a finite subset of $X(\Omega)$, \inserted{$g$ the genus of $X$, $n=\mathrm{card}\,S$.} $\deleted{\unsure{X_0}}\inserted{U}=X\setminus S$, Equivalent conditions

a) $\Pi_1(U)$ nonabelian

b) $\mathrm{Aut}(U)$ finite

c) For $\forall$ \inserted{connected \deleted{of finite type}} reduced scheme $X$ \inserted{of finite type} over $\Omega$, the set of \uline{nonconstant} morphisms from $X$ to $U$ is finite

d) We are \deleted{\unsure{$X$}} in one of the following three cases: $1^\circ)$ \deleted{genus} $g\geq 2$ \quad $2^\circ)$ $g=1$, $n\geq 1$ \quad $3^\circ)$ $g=0$, $n\geq 3$.

\margin{e) (if $\Omega\unsure{\in|\subset}\mathbb C$) The universal covering of $X(\mathbb C)\setminus S(\mathbb C)$ is isomorphic to the half-plane \uline{of} Poincaré.}
\margin{\unsure{?!|??} f) (if $\Omega=\overline{\mathbb Q}$ \unsure{$S\neq\emptyset$}) The universal covering of $X\setminus S=U$ is isomorphic to that of $\mathbb P^1_\Omega\setminus\{0,\unsure{1},\infty\}$.}

\uline{Definition 5.3.} We then say that $(X,S)$ is anabelian. \deleted{If $X$ is a scheme}

As this condition is (by d) for example) invariant under extension of the algebraically closed base field, we extend this definition to the case of a pair $(X,S)$, with $(X,S)$ \deleted{projective smooth of dimension} \inserted{as in (5.1)} (\uline{N.B.} we look separately at the connected components of $X_{\overline K}$\dots).

\margin{19}